\section*{Chapitre \ref{ch:mx:circuit-breakers-and-halts} --- Disjoncteurs et suspensions}

\begin{solution}{exo:mx:circuit-breakers-and-halts:1}
$99.93\times0.9985=99.780$ arrondi à la baisse à 99.78, et $99.93\times1.0015=100.080$ arrondi à la hausse à 100.08.
\end{solution}

\begin{solution}{exo:mx:circuit-breakers-and-halts:2}
15 minutes~: le seuil a été franchi à 9 h 34 min 13 s et les \omterm{def:mx:circuit-breakers-and-halts:pause}{enchères de réouverture} ont commencé à 9 h 49 min 13 s. Une baisse de niveau~2 survenant à 15 h 25 ou après n'aurait pas
suspendu le marché.
\end{solution}

\begin{solution}{exo:mx:circuit-breakers-and-halts:3}
43.9~\% sans mécanisme, 41.1~\% avec le seuil, 29.4~\% avec la \omterm{def:mx:circuit-breakers-and-halts:velocity}{logique de vélocité}, 23.2~\% avec la suspension générale du marché.
\end{solution}

\begin{solution}{exo:mx:circuit-breakers-and-halts:4}
Sa borne inférieure fait office de barrière~: les cours acheteurs situés en dessous sont rejetés et les ordres de vente au marché s'arrêtent à ce niveau, leur reliquat étant annulé. Les ventes qui auraient continué à balayer
davantage le carnet ne sont jamais exécutées, de sorte que le cours ne peut pas passer sous le seuil, alors que celui-ci ne fait que suivre à la baisse le cours moyen de la minute.
\end{solution}

\begin{solution}{exo:mx:circuit-breakers-and-halts:5}
Ses pauses sont brèves, et chaque \omterm{def:mx:circuit-breakers-and-halts:pause}{enchère de réouverture} se heurte à la même vente unilatérale~: la pause interrompt les échanges sans attirer de nouveaux acheteurs, et la baisse
reprend~; pendant ce temps, les ordres de vente envoyés durant la pause attendent la réouverture et alourdissent son déséquilibre.
\end{solution}

\begin{solution}{exo:mx:circuit-breakers-and-halts:6}
Conditionner à la distance à la limite, mesurer la vitesse d'approche ou la probabilité d'atteindre la limite, et comparer avec un groupe témoin
dépourvu du mécanisme (un autre marché, une période antérieure à la règle, une limite qui n'était pas en vigueur)~; distinguer la barrière mécanique (ordres refusés à
la limite) du comportement (ordres avancés dans le temps), par exemple à l'aide de données au niveau des ordres indiquant qui envoie des ordres près de la limite.
\end{solution}

\begin{solution}{exo:mx:circuit-breakers-and-halts:7}
À 3 ticks~: effet d'aimant $0.727-0.694=+0.03$ avec une erreur type d'environ 0.09~; effet de barrière $0.694-0.816=-0.12$ avec environ 0.09. Aucun des deux ne dépasse deux
erreurs types à une distance donnée~; le signe de l'effet de barrière est le même pour dix des onze distances.
\end{solution}

\begin{solution}{exo:mx:circuit-breakers-and-halts:8}
Un moindre nombre de titres atteignant leur limite peut signifier que le seuil bloque les transactions à sa borne (effet de barrière), et non que le marché est plus calme~; le volume et les exécutions refusées
au seuil, ainsi que les cours qui suivent, indiquent si le risque a été réduit ou seulement reporté. Une comparaison avant-après ignore en outre ce qui a pu changer par ailleurs.
\end{solution}

\begin{solution}{pb:mx:circuit-breakers-and-halts:1}
\textbf{1.} Voir les définitions.
\textbf{2.} Des seuils contre les cours qui s'emballent ou sont erronés, des pauses contre les mouvements plus rapides que le retour de la liquidité, des suspensions contre les nouvelles non digérées, des limites pour toute la journée.
\textbf{3.} Par les contrôles du moteur~: R fixe un seuil, P modifie la phase~; la politique vérifie toutes les demi-secondes.
\textbf{4.} 7~\%, 13~\% (15 minutes, avant 15 h 25) et 20~\% (reste de la journée)~; suspensions de niveau~1 les 9, 12, 16 et 18 mars 2020.
\textbf{5.} 40\,000 actions vendues sous forme de soixante \omterm{def:mx:the-limit-order-book:orders}{ordres au marché} sur une minute~; sans mécanisme, un seuil dynamique avec pauses, la \omterm{def:mx:circuit-breakers-and-halts:velocity}{logique de vélocité}, une suspension générale du marché.
\textbf{6.} Baisses de 27.4, 20.0, 24.9 et 18.9 ticks~; 0, 0, 79 et 48 secondes de suspension~; 17\,580, 16\,440, 11\,752 et 9\,279 actions exécutées~; coûts de 12.2, 9.6,
13.2 et 6.6 ticks.
\textbf{7.} Elle interrompt toute négociation au premier niveau et retire du marché le plus grand nombre d'ordres.
\textbf{8.} Une pause après un mouvement d'un nombre fixé de ticks dans un délai fixé~; des ordres peuvent être saisis, modifiés et annulés pendant l'état réservé.
\textbf{9.} Les disjoncteurs peuvent amener les agents à avancer leurs transactions, et donc accroître la variabilité et aggraver les mouvements.
\textbf{10.} Des vendeurs qui, lorsqu'un seuil est appliqué, vendent d'autant plus vite que le meilleur cours acheteur est proche de sa borne inférieure.
\textbf{11.} Pour chaque séance et chaque distance, le premier instant où le cours se trouve à moins de $x$ ticks et la question de savoir si le seuil est atteint dans les 30 secondes~; seuil avec anticipateurs,
seuil seul, absence de seuil.
\textbf{12.} \emph{Résultat nommé}~: à 5 ticks de la borne inférieure, la probabilité de l'atteindre dans les 30 secondes est de 0.66 avec le seuil et
des vendeurs anticipateurs, de 0.63 avec le seuil seul et de 0.75 sans seuil~; l'effet d'aimant simulé vaut en moyenne $+0.03$ sur 3 à 7 ticks (non
distinguable de zéro sur 80 séances), l'effet de barrière du seuil $-0.12$.
\textbf{13.} Il supprime les ordres de vente à sa borne, ce qui l'emporte sur les ventes supplémentaires qu'il suscite.
\textbf{14.} Des cours indicatifs pendant l'enchère, des collars sur le cours de réouverture, des prolongations lorsque le déséquilibre est important, et suffisamment de temps pour que la contrepartie
arrive.
\textbf{15.} Les disjoncteurs londoniens sur actions individuelles ont amélioré la liquidité et réduit la volatilité des titres qui ont continué à être négociés.
\textbf{16.} Elles ont pris fin avant l'arrivée de nouveaux acheteurs.
\textbf{17.} Davantage de traders anticipateurs, ou un seuil qui suspend au lieu de refuser les transactions.
\textbf{18.} Protéger les cours en refusant des transactions~: des baisses plus courtes, moins d'actions vendues.
\textbf{19.} Que les seuils et les suspensions limitent les baisses surtout en refusant des transactions, que l'éventuel effet d'aimant est faible à côté, et qu'une pause n'aide que si
la liquidité peut arriver avant sa fin.
\textbf{20.} Du temps pour que la liquidité revienne, payé par des transactions qui n'ont pas lieu.
\end{solution}

\begin{solution}{iq:mx:circuit-breakers-and-halts:1}
Arrêter l'horloge de l'algorithme, décider si l'on participe à l'\omterm{def:mx:circuit-breakers-and-halts:pause}{enchère de réouverture} (et à quelle limite, compte tenu du cours indicatif), et replanifier le reste
de l'ordre pour après la réouverture~; ne jamais envoyer d'\omterm{def:mx:the-limit-order-book:orders}{ordres au marché} pendant l'enchère.
\iqlookfor{Participation à la réouverture~; replanification.}
\end{solution}

\begin{solution}{iq:mx:circuit-breakers-and-halts:2}
L'accélération des cours vers une limite, parce que les traders se précipitent avant qu'elle ne devienne contraignante~; la tester en conditionnant à la distance à la limite et en comparant avec un
marché ou une période qui en sont dépourvus, en séparant la barrière du comportement.
\iqlookfor{Conditionnement à la distance~; groupe témoin~; mécanique contre comportement.}
\end{solution}

\begin{solution}{iq:mx:circuit-breakers-and-halts:3}
Les actions au comptant et les options s'arrêtent, mais les contrats à terme peuvent continuer à être négociés dans leurs propres limites, de sorte que les couvertures ne peuvent être ajustées que là où les marchés sont ouverts~; les limites de
risque doivent utiliser les derniers cours avec prudence, et les appels de marge peuvent survenir avant la réouverture.
\iqlookfor{Différences entre marchés~; cours périmés.}
\end{solution}

\begin{solution}{iq:mx:circuit-breakers-and-halts:4}
Une suspension~: arrêter et attendre la réouverture, annuler ou conserver les ordres selon la politique~; une pause~: participer à l'\omterm{def:mx:circuit-breakers-and-halts:pause}{enchère de réouverture} ou attendre~; un rejet par le seuil~:
recoter dans le seuil ou attendre, ne jamais réessayer aveuglément~; dans tous les cas, consigner et alerter.
\iqlookfor{Gestion selon l'état~; pas de nouvelle tentative aveugle.}
\end{solution}

\begin{solution}{iq:mx:circuit-breakers-and-halts:5}
La \omterm{def:mx:circuit-breakers-and-halts:velocity}{logique de vélocité} réagit à la vitesse (ticks par seconde) avec une pause de quelques secondes~; le Limit Up--Limit Down (LULD) réagit à la distance à une référence mobile, avec un
état de limite et une pause de cinq minutes~: l'une porte sur les mouvements rapides, l'autre sur les mouvements amples.
\iqlookfor{Vitesse face à la distance~; durée des pauses.}
\end{solution}

\begin{solution}{iq:mx:circuit-breakers-and-halts:6}
Des données sur la fréquence à laquelle les seuils deviennent contraignants et les pauses se déclenchent, sur la poursuite des cours après les réouvertures et sur les effets de contagion vers d'autres titres~; s'attendre à moins de déclenchements
et à moins de volumes refusés, mais à des mouvements plus amples à l'intérieur des seuils et à moins de temps pour que la liquidité revienne.
\iqlookfor{Compromis~; mesures à surveiller.}
\end{solution}
